\section{总结与下一步工作}
\label{sec:conclusion}

本任务已经形成从数据解析、基础潮流到临界稳定分析的工作流。NR 具有近解区域的快速收敛能力，FDLF 具有较低的单轮计算成本，但对高 $R/X$ 和强耦合工况更敏感。接近稳定极限后，算法是否返回结果不能单独证明物理解存在，必须结合最终残差、无功限值、雅可比奇异性和连续潮流分支综合判断。

\subsection{当前结论}

\begin{enumerate}
  \item 常规工况下 NR 与 FDLF 均能可靠求解，前者迭代少，后者单轮成本低；
  \item 高 $R/X$、重载及鼻点附近工况会显著削弱解耦近似和局部求根的有效性；
  \item 全局化方法的作用是控制迭代过程，不能改变原潮流方程是否存在实数根；
  \item CPF 与奇异性指标应作为最大负荷能力判断的主要依据。
\end{enumerate}

\subsection{当前不足}

目前 CPF 参考鼻点与定点算法收敛边界尚未在完全统一的控制方程和无功限值切换口径下逐点复核；规模分析也仍基于当前 Python 与矩阵实现，尚未区分算法复杂度和代码实现开销。上述两点是下一阶段需要优先消除的不确定性。

\subsection{下一步工作}

后续工作建议如下：
\begin{enumerate}
  \item 统一不同求解器的负荷参数化、PV--PQ 转换与残差口径；
  \item 增加雅可比最小奇异值、条件数和临界特征向量监测；
  \item 引入稀疏矩阵存储与稀疏 LU/QR 分解；
  \item 对多种 CPF 参数化开展统一的精度、成本和跨鼻点能力比较；
  \item 保存配置、代码版本、原始数据和图像来源，提高结果可复现性。
\end{enumerate}
